\appendix
\section{DigitalMindExplorer}
\label{app:framework}

\textsc{DigitalMindExplorer} is the runtime the probe is built on. An experiment
about a mind needs the mind to be inspectable. An agent framework built for task
performance hides exactly the parts an experimenter wants to hold, so we wrote
one that exposes them. The package has no required dependencies, and a provider
is imported only when a model from it is asked for.

\paragraph{What a mind is made of.} A mind holds modules. A module has a queue,
a set of declared output channels, and a turn. A turn is two steps: one
\texttt{on\_input} for each message that arrived, in order, then one
\texttt{on\_process}. Splitting them means a module absorbs everything that
reached it before it acts. An agent is a module with a model, a transcript, and
instrumented calls. Subject and ego are agents, and a stage is whatever the
experiment needs it to be.

\paragraph{Wiring.} Registering is the only wiring verb. A producer registers a
consumer onto one of its declared channels, optionally renaming the channel on
arrival, so neither module needs to know the other's vocabulary. Registering
against a channel that was never declared is an error, so a typo fails when the
mind is assembled rather than by silently dropping messages. The mind itself
holds no routing table. It offers one shorthand, \texttt{intercept}, which places
stages between the subject and the ego, and that shorthand is exactly a sequence
of register calls.

\paragraph{Determinism.} The scheduler is the two-phase tick described in
Section~3. Modules take their turns concurrently, and every emission is staged
into a private outbox that is delivered at the end of the tick. Two agents can
therefore think at the same moment without either one observing how far the other
got. Delivery order is fixed by registration order. An external input runs to
quiescence, so the mind has finished thinking when the call returns.

\paragraph{Instrumentation.} Every message, model call, memory write, and turn
boundary is an event with a timestamp, a tick, and the module that produced it.
Events are fanned out to sinks. A run writes one trace of everything in order and
one file per module, so a single part can be read on its own. Text payloads are
stored in full rather than summarised, which is what makes a trace a recording
rather than a description. Attaching an extra sink is one line and changes
nothing else, which is how the browser interface receives the same events live.

\paragraph{Swapping models.} A model is named by a provider string, so moving a
part from a local open-weight model to a hosted one changes one string. Wrapping
a whole mind in a recording is one argument at construction. Every call is then
taped once and replayed afterwards without a model, which is what makes a
published run checkable by someone who cannot run the original.

\paragraph{Reference architectures.} Three minds ship with the runtime, one per
file. \emph{Plain} is a subject on its own. \emph{Interceptor} places an editor
between the subject and the ego, which is the shape the probe in this paper uses.
\emph{Bicameral} places an inner voice that inserts an unbidden utterance into
the window, and attaches a workspace that observes every channel. Each declares a
title, the shape of its pipeline, and which extra models it takes, so a script or
the interface can list and build any of them without knowing the details.

\paragraph{Watching a mind run.} A browser interface opens on a chooser: the
minds available to run, and the runs already recorded. Talking to a mind shows
each part's context window side by side, in pipeline order, with each window
compared against the one it derives from. The comparison is a diff of two message
lists, so an edit, an insertion, and a deletion are all reported without the
interface knowing what any stage does. Selecting a recorded run replays it, one
event or one tick at a time, rebuilt from the trace alone.
